\documentclass{jsarticle}
\usepackage{amsmath}
\usepackage{amssymb}
\begin{document}
\title{坂本龍一先生のアルゴリズムを、見ると}
\author{Masaaki Yamaguchi}
\date{}
\maketitle



         綺麗な和音の連鎖の基本は、B♭とE♭の音階のとりうる、
    (A♯ )(G♯ )(F♯ )(E)(D)(C♯) 、(D♯ )(C♯)(B)(A♯)(G♯)(F♯)
    のエントロピー値をとる、音階をずらす。本当は、後にずらす。
    バイオリンの女医に教えてもらえた。くやしいが、エロ雑誌で同じ人を見ている。
    この女性は、フルメタとは違う。すごく目も綺麗で、端正こもっていた。
    東洋の神秘です。
    初めは、私も単音かなとおもっていたが、この女医は、さもなく、ビアノの
    鍵盤をずらして見せて、音階を一瞬でみせていた。石川先生の助言で
    家に帰って、最近になって、やってみて、なるほど！と思った。
    でも、単音でも、実際の不協和音として、単音の和音として、
    これが、メロディーに分散に入れるのが、難しい。
    これも不協和音と確実に専門家は知っている。これをまべなく、不協和音
    の和音よりも、単独で入れるのは、至難の業である。ハイドが造作もなく
    やっている。楽譜をみせてもらえて、できるか！とおもった。
    あのときは、歌ってみたが、真似ているからできていた。タバコをすっている
    ので、シャウトが心地よく、独特の歌声である。パンクロックがすごく
    好きらしい。キリストの曲もすごく作っている。
    河村隆一は、タバコを吸わないが、ハイドは、アナウンサーのキャスターと
    結婚して、この人は、家庭的なモデルと結婚した。
    両方の歌声が好きだ。河村隆一も、シャウトを使うが、独特のG.とRoserが
    ある。Slaveもすごかった。ヴォーカルによって、タバコで声を変えたり、
    タツゥーで声を根本からかえる人もいる。
    クラシックから入る人もいる。
    河村隆一がその人だ。
    話は、ハイドのほうが好きだ。　


         音階の仕組みを知っていると、楽に、作曲にこの単音の不協和音
    を入れることが、さもなく、直感で鍵盤を見るとわかる。
    苫米地博士は、確実に知っている。バッハのフーガになる。
    石川先生も知っている。１２音と１２音階は違う。
    これを基本に、C♯　と　F♯　G♯　のとりうる音階を
    合わすと、F♯とG♯の、G♭　A♭　B♭、A♭　B♭　とB♭、E♭の音階も、
    これに合わすと、GとAとBの単音をきれいにすると、Gのときは、
    A♯を重ねると、全音階の不協和音がきれいになる。Aのときは、
    F♯を重ねる。あとは、２音階のときは、４オクターブが出来て、
    １音階のときは、３オクターブまでカバーできる。
    C♯は、エナジー・フローと鉄道員が基本の綺麗さであり、
    B♭とE♭は、ラストエンペラーが基本であり、Forbidden Colourも基本である。
    G♭のとりうる音階は、Fがない。F♯のとりうる音階はGがない。
    バッハのフーガは、リトル・ブッダが基本である。音階は、B♭とE♭とC♯である。
    私の無作為の作曲は、全音階の５オクターブをカバーできるらしい。
    あの人は、発見・発明と無作為の音階でB♯でAとして、不協和音では、２パターンの
    音階で４オクターブがカバーできるのを、もっと上の、共鳴なのに、不協和音
    としてのほめことばであった。この共鳴は、不協和音として入れるべきとおもう。
    坂本龍一先生がすごいのは、不協和音をとる音階で曲を見ると、
    作曲の基本がわかるということである。この人は、３秒から４秒聞くと、
    その即興曲の音階がどの音階を使っているかを一瞬でわかる能力が本当である。
    この人を参考に、楽譜を速読しようとすると見えなくなるので、
    私は作曲に向かい、作曲の方法が分かると、楽譜を風景とみると、
    HANONを練習していると、引けるとわかった。HANONは、言葉の文を
    音の文と同等に速読するのに、うってつけだった。私が、作曲できたので
    言えている。不協和音は、難しすぎて、避けていたが、作曲の方法がわかると、
    自分でも、できた。音は使っていない。
    単音をいくらやっても、和音のとりうる、きれいに聞こえる範囲が
    わからない。作曲の音の構成を知ると、どの不協和音で、長調の音を、
    どの変調と嬰調を入れると、不協和音の音階の、長調である一番むずかしいGとAの
    組み合わせが、この理論の全不協和音の音階の中に、GとAを入れることが
    できるのが直感で鍵盤を見て分かるということである。
    百聞は一見にしかずである。
    綺麗に聞こえても、メロディーとしてと、伴奏としてのきれいさは別である。
    このメロディーと伴奏のきれいさのパターンは、坂本龍一先生から、
    学んだ。姉たちのピアノを引いているきれいさからも学んだ。



       一音階では、３オクターブが制限としてあり、これを超えると、魔法の
    メロディーでないと、きれいに聞こえない。
    二音階では、４から５オクターブの範囲が使える。この二音階は、２オクターブ
    と同じ意味ではないのは、今までの説明でわかるとおもう。
    ２オクターブは、１オクターブのまたその上の１オクターブであるが。

    私のB♯は、Cのハ長調では、説明がつかない、不協和音を奏でるので、
    本当に、B♯のAらしい。直感で弾いて、５オクターブ以上を１種音に
    出来ている。Cのハ長調は、１オクターブの繰り返す、上に下に同じ
    メロディーになる音階である。作曲する人は知っているらしい。
    私のB♯は、５オクターブ以上を単一メロディーにできる。
    あの人は、初めは、ト音記号でのC♯をへ音記号でのB♭になっていると
    初め入っていた。坂本先生は、ト音記号をへ音記号のCで循環になっていると
    いっていたが、両者とも、実際に聞いてみて、考えを改めてくれていた。　


       二音階で、４から５オクターブとは、一音階で、１オクターブとはちがう。
    パターンも１オクターブ、２オクターブとも違う。
    一音階で、３オクターブは、１オクターブごとの繰り返しではない。
    ３オクターブもろもろの分散範囲で、曲のメロディーや旋律を形づくる。
    ３オクターブで一音種を作ることである。
    パターンが一種音であるが、言いたいことは、分かると思う。
    ２オクターブは、１オクターブのまたその上の１オクターブであるが。
    １オクターブは、白鍵盤は７音であるが。
    １オクターブは、音では１２音である。
    １２音階は、これとは別。
    二音階は、二和音でもない。三和音でもない。不協和音であるが、和音ではない。
    和音で、二音階をとるのは、長調同士の組み合わせもあり、
    でも、この長調の中にも不協和音があるが。


    ハーバードは、これをさもなくやりこなすらしい。

\end{document}
